\title{xLSTM: Extended Long Short-Term Memory}

\begin{abstract}
The introduction of the constant error carousel and gating in the 1990s established the foundational mechanisms of the Long Short-Term Memory (LSTM) architecture. LSTMs have since become pivotal in the evolution of deep learning, playing a major role in many milestone achievements, notably in forming the groundwork for the first Large Language Models (LLMs). Despite these accomplishments, the emergence of Transformer architectures—characterized by their inherently parallelizable self-attention mechanism—has set a new standard, surpassing LSTMs in scalability and effectiveness. This prompts us to revisit the question: If we scale LSTMs to billions of parameters and incorporate advancements from the latest LLMs while addressing longstanding weaknesses of LSTMs, what level of performance can we achieve in language modeling? To this end, we propose exponential gating enhanced by tailored normalization and stabilization strategies. Furthermore, we present restructured LSTM memories, resulting in: (i) the sLSTM, which employs a scalar memory, scalar updates, and novel mixing operations; and (ii) the mLSTM, which enables full parallelization through a matrix-based memory and a covariance-driven update mechanism. By incorporating these improved LSTM variants within residual block frameworks, we construct xLSTM blocks, which can be stacked to build xLSTM architectures. These innovations—exponential gating and redefined memory structures—substantially improve the capabilities of xLSTM, enabling them to compete effectively with state-of-the-art Transformers and State Space Models in both performance and scaling behavior.\\
Code available at: \url{https://github.com/NX-AI/xlstm}
\end{abstract}

\section{Introduction}
\label{sec:intro}

The concepts underlying Long Short-Term Memory (LSTM)~\citep{Hochreiter:91,Hochreiter:97nips,Hochreiter:97}, such as gating mechanisms and the constant error carousel, were developed specifically to address the vanishing gradient issue encountered in recurrent neural networks~\citep{Hochreiter:91,Hochreiter:00book}:

\begin{align}
\label{eq:lstm_idea}
  c_t \ = \ f_t \ c_{t-1} \ + \ 
          i_t \  z_t \ , \quad  
          h_t \ = \ o_t \ \psi({c_t}) \ . 
\end{align}

The cell state $c_{t-1}$ (green) undergoes an additive update known as the constant error carousel, incorporating the cell inputs $z_t$ and being regulated by sigmoid gates (blue). This update process is governed by the input gate $i_t$ and the forget gate $f_t$, whereas the output gate $o_t$ determines the final output of the memory cell, specifically the hidden state $h_t$. Prior to the output gating, the cell state is compressed or normalized through $\psi$, after which it forms the hidden state.

LSTMs have demonstrated significant success across a range of fields \citep{Hochreiter:01,Hochreiter:07,Schmidhuber:15}, dominating the landscape of text generation prior to the rise of Transformers in 2017~\citep{Vaswani:17}. Their efficacy has been shown through application to numerous sequence-based problems, such as automated text production~\citep{Graves:13,Karpathy:15blog}, handwriting synthesis~\citep{Graves:13}, sequence-to-sequence translation tasks~\citep{Sutskever:14nips}, program evaluation~\citep{Zaremba:14arxiv}, caption generation for images~\citep{Karpathy:15,Hossain:19}, source code synthesis~\citep{Karpathy:15blog}, hydrological time series modeling for rainfall and runoff~\citep{Kratzert:18,Kratzert:19}, and real-time hydrological forecasting, including flooding alerts~\citep{Nearing:24}. In reinforcement learning domains, LSTMs continue to be leading architectures for sequential modeling, as exemplified by models like AlphaStar in StarCraft II~\citep{Vinyals:17short}, OpenAI Five for Dota 2~\citep{OpenAI:19}, and nuclear fusion magnetic controller models~\citep{Degrave:22short}. Their remarkable ability to capture abstract representations—by extracting and retaining semantic features within their cell states~\citep{Karpathy:15blog}—has been observed through the discovery of specialized neurons responsible for numbers and syntax~\citep{Lakretz:19}, linguistic processing~\citep{Bau:19}, and sentiment detection~\citep{Radford:17arxiv}. Even now, LSTMs continue to underpin pivotal real-world systems~\citep{Degrave:22short,Nearing:24}, illustrating their enduring utility and robustness.

Although LSTMs have achieved notable success, they face three significant drawbacks: (i) Inflexibility in modifying prior memory decisions. This issue can be illustrated through the {\it Nearest Neighbor Search} problem (refer also to Appendix~\ref{sec:appNearestNeighborSearch}). Here, given a reference vector, the model must identify the most similar vector within a sequence and produce its associated value at the sequence's conclusion. As presented in the left panel of Figure~\ref{fig:lstmProblems}, LSTM displays a higher mean squared error in this context, as it cannot efficiently update the stored value when encountering a vector of greater similarity. In contrast, our xLSTM addresses this shortcoming with its exponential gating mechanism. (ii) Constrained storage capability, which forces information to be compressed into one-dimensional cell states. This constraint becomes apparent in the {\it Rare Token Prediction} setting. As depicted in the right panel of Figure~\ref{fig:lstmProblems}, LSTM exhibits higher perplexity in predicting rare tokens in the Wikitext-103 dataset~\citep{Merity:17}, due to its limited storage scope. Our xLSTM alleviates this by utilizing a matrix-based memory. (iii) Absence of support for parallel computation, stemming from the recurrence links between consecutive hidden states that necessitate step-by-step sequential operations.

The shortcomings inherent in LSTM have led to the rise of Transformers~\citep{Vaswani:17} within the field of language modeling. It is thus pertinent to ask what results in language modeling might be attainable if we address these weaknesses and scale up LSTMs to match the magnitude of today’s Large Language Models.

\section{Extended Long Short-Term Memory}
\label{sec:xLSTM}

In response to the limitations of LSTM, Extended Long Short-Term Memory (xLSTM) implements two principal innovations beyond the core LSTM formulation in Equation~\eqref{eq:lstm_idea}. These innovations — exponential gating and alternative memory architectures — expand the LSTM framework into two distinct variants: (i) the sLSTM (detailed in Section~\ref{sec:sLSTM}), characterized by a scalar memory, scalar updates, and integrated memory mixing; and (ii) the mLSTM (see Section~\ref{sec:mLSTM}), which utilizes matrix-based memory representations and adopts a covariance (outer product) update mechanism that is amenable to full parallel computation. Both sLSTM and mLSTM utilize exponential gating mechanisms to improve upon the traditional LSTM. To facilitate parallelism, the mLSTM forgoes memory mixing—that is, there are no recurrent hidden-to-hidden connections. Both mLSTM and sLSTM architectures can be generalized to handle multiple memory cells. In this setup, sLSTM allows for memory mixing between different cells, while additionally supporting multi-head designs where mixing only occurs among cells within the same head and not across heads. The introduction of heads, together with exponential gating, offers a new paradigm for memory mixing within sLSTM. For mLSTM, the concepts of multi-headedness and multiple cells are functionally synonymous.

The incorporation of these novel LSTM variants within residual block modules produces what we refer to as xLSTM blocks (refer to Section~\ref{sec:xLSTMblock}). Constructing neural architectures by residually stacking these xLSTM blocks gives rise to xLSTM architectures, as further detailed in Section~\ref{sec:xLSTMarchitecure}. An overview of the xLSTM architecture and its constituent elements can be found in Figure~\ref{fig:xlstm_sketch}.

\subsection{Review of the Long Short-Term Memory}
\label{sec:orgLSTM}

The initial concept for LSTM~\citep{Hochreiter:91,Hochreiter:97nips,Hochreiter:97} featured a scalar memory cell as its core computational and storage component, designed to circumvent the problem of vanishing gradients~\citep{Hochreiter:91,Hochreiter:00book} via the constant error carousel, which operates through updates to the cell state. The memory cell architecture employs three gating mechanisms—input, output, and forget gates—with the latter being subsequently introduced by \citet{Gers:00}. The mechanisms that govern the LSTM memory cell's updates at time step $t$ are as follows:

\begin{align}
c_t \ &= \ f_t \, c_{t-1} \ + \ i_t \, z_t &  & &\text{cell state} \\
h_t \ &= \ o_t \, \tilde{h}_t \ , & \tilde{h}_t \ &= \ \psi \left( c_t\right)
  &\text{hidden state} \\
z_t \ &= \ \varphi \left( \tilde{z}_t \right) \ , 
  &\tilde{z}_t \ &=  \ \mathbf{w}^\top_{z} \, \mathbf{x}_t \ + \
  r_{z} \, h_{t-1} \ + \  b_{z} \ \
  &\text{cell input} \\
i_t \ &= \ \sigma \left( \tilde{i}_t  \right) \ , 
  &\tilde{i}_t \ &= \ \mathbf{w}^\top_{i} \, \mathbf{x}_t \ + \
  r_{i} \, h_{t-1} \ + \  b_{i} \ \
  &\text{input gate} \\
f_t \ &= \ \sigma \left(  \tilde{f}_t \right) \ , 
  &\tilde{f}_t \ &= \ \mathbf{w}^\top_{f} \, \mathbf{x}_t  \ + \
  r_{f} \, h_{t-1} \ + \  b_{f} \ \
  &\text{forget gate} \\
o_t \ &= \ \sigma \left( \tilde{o}_t \right) \ , 
  &\tilde{o}_t  \ &= \ \mathbf{w}^\top_{o} \, \mathbf{x}_t \ + \
  r_{o} \, h_{t-1} \ + \  b_{o} \ \
  &\text{output gate} 
\end{align}

The vectors $\mathbf{w}_{z}$, $\mathbf{w}_{i}$, $\mathbf{w}_{f}$, and $\mathbf{w}_{o}$ denote the input weights mapping from the input vector $\mathbf{x}_t$ to the cell input, input gate, forget gate, and output gate, accordingly. The recurrent weight parameters $r_{z}$, $r_{i}$, $r_{f}$, and $r_{o}$ link the preceding hidden state $h_{t-1}$ to the cell input and each respective gate (input, forget, and output). The terms $b_{z}$, $b_{i}$, $b_{f}$, and $b_{o}$ represent biases for their respective components. Activation functions $\varphi$ and $\psi$ act on the cell input and hidden state, typically instantiated as the $\tanh$ function. The activation function $\psi$ serves to restrict the range of the cell state, preventing it from diverging. All gate computations are performed using the sigmoid function, $\sigma \left( x \right) = 1/(1+\exp(-x))$. Subsequent advancements bundled several scalar memory units into the vector $\mathbf{c}_t \in \mathbb{R}^{d}$, facilitating the use of recurrent weight matrices $\mathbf{R} \in \mathbb{R}^{d \times d}$ to interconnect the outputs of distinct memory units~\citep{Greff:15}. Further elaboration is provided in Appendix~\ref{sec:apporgLSTM}. Investigations using ablation analysis have demonstrated the essential role of every individual component within the memory cell~\citep{Greff:15}.

\subsection{sLSTM}
\label{sec:sLSTM}

In order to enhance LSTMs with the capacity to adjust their memory retention choices, we implement exponential gates (depicted in red) in conjunction with methods for normalization and stabilization. Specifically, we enable the input and forget gates to utilize exponential activation functions. For normalization purposes, we define a normalizer state, which accumulates the sum of products between the input gate and all subsequent forget gates.

The scalar sLSTM forward pass is:

\begin{align}
\label{eq:slstmforward}
c_t \ &= \  f_t \ c_{t-1} \ + \ i_t \ z_t & & 
 &\text{cell state} \\
n_t \ &= \  f_t \ n_{t-1} \ + \ i_t  & & 
 &\text{normalizer state} \\
h_t \ &= \ o_t \ \tilde{h}_t \ ,
  &\tilde{h}_t \ &= \ c_t / n_t
&\text{hidden state} \\
z_t \ &= \ \varphi \left( \tilde{z}_t \right) \ , 
  &\tilde{z}_t \ &=  \ \mathbf{w}^\top_{z} \ \mathbf{x}_t \ + \
  r_{z}  h_{t-1} \ + \  b_{z}
  &\text{cell input} \\
i_t \ &= \ \!\exp \left( \tilde{i}_t  \right) \ , 
  &\tilde{i}_t \ &= \ \mathbf{w}^\top_{i} \ \mathbf{x}_t \ + \
  r_{i}  \ h_{t-1} \ + \  b_{i}
  &\text{input gate} \\
f_t \ &= \ \sigma \left(  \tilde{f}_t \right) \ \text{OR} \
\!\exp \left(  \tilde{f}_t \right) \ ,
  &\tilde{f}_t \ &= \ \mathbf{w}^\top_{f} \ \mathbf{x}_t  \ + \
  r_{f}  \ h_{t-1} \ + \  b_{f}
  &\text{forget gate} \\
o_t \ &= \ \sigma \left( \tilde{o}_t \right) \ , 
  &\tilde{o}_t  \ &= \ \mathbf{w}^\top_{o} \ \mathbf{x}_t \ + \
  r_{o}  \ h_{t-1} \ + \  b_{o}
  &\text{output gate} 
\end{align}

We adapt the traditional LSTM gate mechanisms—incorporating gating based on input and/or hidden states as well as a bias component—into our proposed architectures. The use of exponential activation functions can result in excessively large outputs and potential numerical overflow. To address this, we employ a stabilization approach for the gates through the introduction of an auxiliary state, denoted as \mbox{$m_t$~\citep{Milakov:18arxiv}}:

\begin{align}
\label{eq:slstmstabil}
m_t \ &= \ \max \left( \log ( f_t ) + m_{t-1} , \log ( i_t ) \right) &\text{stabilizer state} \\
i'_t \ &= \ \!\exp \left( \log \left ( i_t \right) - m_t \right) \ = \ \!\exp \left( \tilde{i}_t - m_t \right) \ \, 
  &\text{stabil. input gate} \\
f'_t \ &= \ \!\exp \left( \log \left( f_t \right) + m_{t-1} - m_t \right) \ \, 
  &\text{stabil. forget gate}
\end{align}

\noindent

The equations above specify a stabilized variant of gate computations in recurrent models. Equation \eqref{eq:slstmstabil} defines the stabilizer state, $m_t$, as the greater of $\log(f_t) + m_{t-1}$ and $\log(i_t)$, thereby controlling scaling in the cell. The stabilized input gate $i'_t$ is calculated by exponentiating the difference between the logarithm of $i_t$ and the stabilizer $m_t$ (or $\tilde{i}_t - m_t$ equivalently), ensuring numerical stability. Likewise, the stabilized forget gate $f'_t$ is evaluated by exponentiating the sum of $\log(f_t)$ and $m_{t-1}$ minus $m_t$, thereby normalizing the gate's magnitude.

As demonstrated in Appendix~\ref{sec:appsLSTM}, substituting $f_t$ with $f'_t$ and $i_t$ with $i'_t$ during the forward computation leaves both the network’s overall output and the gradients of the loss function with regard to the parameters unaffected.

\paragraph{New Memory Mixing.} Similar to the conventional LSTM, sLSTM is capable of supporting multiple memory cells (refer to Appendix~\ref{sec:appsLSTM}). The use of multiple memory cells facilitates the process of memory mixing by employing recurrent connections $\mathbf{R}_{\mathbf{z}}$, $\mathbf{R}_{\mathbf{i}}$, $\mathbf{R}_{\mathbf{f}}$, and $\mathbf{R}_{\mathbf{o}}$ that link the hidden state vector $\mathbf{h}$ to the memory cell input $\mathbf{z}$ and the gates $\mathbf{i}$, $\mathbf{f}$, and $\mathbf{o}$, respectively. What distinguishes the current form of memory mixing is the role of exponential gating. In the reconfigured sLSTM, there can be several heads, each performing memory mixing internally, but interactions do not occur between different heads. The combination of these heads and exponential gating in sLSTM introduces a novel mechanism for memory mixing.

\subsection{mLSTM}
\label{sec:mLSTM}

In order to improve the storage capability of LSTMs, we generalize the LSTM memory cell from being a scalar $c \in \mathbb{R}$ to a full matrix $\mathbf{C} \in \mathbb{R}^{d \times d}$. This adjustment means that memory retrieval now employs matrix multiplication. Specifically, at each time step $t$, a pair of vectors are to be stored: the key $\mathbf{k}_t \in \mathbb{R}^d$ and the value $\mathbf{v}_t \in \mathbb{R}^d$, adopting terminology from the Transformer architecture. At some subsequent time $t + \tau$, retrieval of $\mathbf{v}_t$ is carried out using a query vector $\mathbf{q}_{t+\tau} \in \mathbb{R}^d$. This mechanism parallels the framework known as Bidirectional Associative Memories (BAMs) \citep{Kohonen:72,Anderson:72,Nakano:72,Anderson:77}.

The covariance update rule~\citep{Sejnowski:77,Dayan:91} for storing a key-value pair is

\begin{align}
\mathbf{C}_t \ &= \ \mathbf{C}_{t-1} \ + \ \mathbf{v}_{t} \ \mathbf{k}_t^\top \ .
\end{align}

We presuppose the application of layer normalization before projecting the inputs to keys and values, resulting in these vectors having zero mean. The rule for updating the covariance matrix has been demonstrated to be optimal~\citep{Dayan:91} in achieving maximal separability among binary vectors during retrieval—a property equivalent to attaining the highest possible signal-to-noise ratio. Restricting retrieval to interactions between pairs and allowing for quadratic computational costs, as is characteristic of attention mechanisms~\citep{Krotov:16,Krotov:17,Ramsauer:21}, can further enhance this separability. Additionally, the covariance update mechanism aligns with the principles of Fast Weight Programmers~\citep{Schmidhuber:92ncfastweights,Schlag:21}, which were subsequently extended by incorporating both a fixed decay coefficient applied to $\mathbf{C}_{t-1}$ and a fixed learning rate scaling the outer product $\mathbf{v}_{t} \mathbf{k}_t ^\top$~\citep{Ba:16}. Following this perspective, we embed the covariance update rule within the LSTM architecture, interpreting the forget gate as the decay rate, the input gate as the learning rate, and the output gate as modulating the scale of the retrieved vector.

Within this matrix memory framework, the normalizer state is constructed by aggregating the key vectors, each scaled by the corresponding input gate as well as the product of subsequent forget gates. The normalizer thus captures the cumulative influence of the gates over time. To address scenarios where the dot product between a query vector and the normalizer state approaches zero, we adopt the approach proposed previously~\citep{Sun:23arxiv}: the absolute value of the dot product is used and is clipped from below at a set threshold (commonly 1.0). The resulting mLSTM computation is as follows:

\begin{align}
\label{eq:mlstm_recurrent_begin}
\mathbf{C}_t \ &= \ f_t \ \mathbf{C}_{t-1} \ + \ 
  i_t \ \mathbf{v}_t \ \mathbf{k}_t^\top & & &\text{cell state} \\
\mathbf{n}_t \ &= \ f_t \ \mathbf{n}_{t-1} \ + \ 
  i_t \ \mathbf{k}_t & & &\text{normalizer state} \\
\mathbf{h}_t \ &= \ \mathbf{o}_t \ \odot \ \tilde{\mathbf{h}}_t
\ , \qquad \qquad 
\tilde{\mathbf{h}}_t \ = \  \mathbf{C}_t \mathbf{q}_t \ / \ 
  \max \left\{ |\mathbf{n}_t^\top \mathbf{q}_t|, 1 \right\}
   &&&\text{hidden state} \\
\mathbf{q}_t \ &= \ \mathbf{W}_q \ \mathbf{x}_t \ + \ \mathbf{b}_q & & &\text{query input} \\
\mathbf{k}_t \ &= \ \frac{1}{\sqrt{d}} \mathbf{W}_k \ \mathbf{x}_t \ + \ \mathbf{b}_k & & &\text{key input} \\
\mathbf{v}_t \ &= \ \mathbf{W}_v \ \mathbf{x}_t \ + \ \mathbf{b}_v & & &\text{value input} \\
i_t \ &= \ \exp \left( \tilde{i}_t \right) \ , 
 \qquad \qquad \quad
  \tilde{i}_t \ = \ \mathbf{w}^\top_{i} \ \mathbf{x}_t \ + \  b_{i} 
  &&&\text{input gate} \\
f_t \ &= \ \sigma \! \left( \tilde{f}_t \right) \ \text{OR} \ \exp \! \left( \tilde{f}_t \right) , \ \
\tilde{f}_t \ = \ \mathbf{w}^\top_{f} \ \mathbf{x}_t  \ + \
  b_{f}
  &&&\text{forget gate} \\
\label{eq:mlstm_recurrent_end}
\mathbf{o}_t \ &= \ \sigma \left( \tilde{\mathbf{o}}_t \right) \ , \qquad \qquad \quad
  \tilde{\mathbf{o}}_t  \ = \ \mathbf{W}_{\mathbf{o}} \ \mathbf{x}_t \ + \
  \mathbf{b}_{\mathbf{o}}
  &&&\text{output gate}
\end{align}

The mLSTM architecture is capable of supporting several memory cells, akin to the standard LSTM structure. In the context of mLSTM, utilizing multiple heads yields the same effect as employing multiple cells, due to the absence of memory mixing. To ensure stability for the exponential gating mechanisms within mLSTM, we apply the identical stabilization methods used for sLSTM (refer to Equation~\ref{eq:slstmstabil}). Given that memory mixing is not present in mLSTM, its recurrence relation permits a reformulation into a parallelizable variant. Additional explanations can be found in Appendix~\ref{sec:appmLSTM}.

\subsection{xLSTM Architecture}
\label{sec:xLSTMarch}

\paragraph{xLSTM Blocks.}
\label{sec:xLSTMblock}
An xLSTM block is designed to capture a non-linear representation of past sequences within a high-dimensional space, thereby enabling clearer distinction between varying histories or contexts. Achieving effective historical separation is essential for accurate prediction of subsequent elements in a sequence, such as the next token. Our approach is guided by Cover’s Theorem~\citep{Cover:65}, which asserts that data points are more likely to be linearly separable when non-linearly projected into a space of greater dimensionality than in the original feature space. 

We explore two variations of residual block architectures: (i) The residual block employing post up-projection, analogous to the architecture in Transformers, first applies a non-linear summary of the past within the original dimensionality, followed by a linear transformation that projects to a high-dimensional space, applies a non-linear activation, and then projects the output linearly back to the original space; this is illustrated in the left panel of Figure~\ref{fig:Backbones} and the third column of Figure~\ref{fig:xlstm_sketch}. A comprehensive depiction is available in Figure~\ref{fig:appsLSTM_detailed} in the appendix. (ii) The second approach is a residual block with pre up-projection, similar to those found in State Space Models, where the transition to high-dimensional space is achieved first through a linear mapping.

The past is first non-linearly encoded within a high-dimensional space, after which it undergoes a linear transformation to return to the original feature space. When an xLSTM block incorporates an sLSTM, the up-projection operation is applied after the block. Conversely, for xLSTM blocks utilizing an mLSTM, the up-projection precedes the block due to the increased memory capacity available in the higher-dimensional context. Further elaboration is available in the left panel of Figure~\ref{fig:Backbones}, the third column of Figure~\ref{fig:xlstm_sketch}, and Figure~\ref{fig:appsLSTM_detailed} found in the appendix.

\paragraph{xLSTM Architecture. }
\label{sec:xLSTMarchitecure}
The xLSTM architecture is developed through the residual stacking of modular components~\citep{Srivastava:15,He:16}. Our design is based on widely adopted pre-LayerNorm~\citep{Ba:16b} residual backbones, mirroring the structure found in current Large Language Models. This configuration is illustrated in the final column of Figure~\ref{fig:xlstm_sketch}.

\subsection{Memory and Speed Considerations}

In contrast to Transformers, xLSTM networks exhibit linear computational costs and maintain fixed memory requirements regardless of the input sequence length. Due to their compressive memory capabilities, xLSTMs are particularly advantageous for deployment in industrial scenarios and edge device implementations.

The memory utilized by mLSTM is parameter-free; however, it incurs substantial computational overhead due to the necessity of maintaining and updating a $d \times d$ matrix. This results in a balance between available memory capacity and computational demands. Nonetheless, since these operations are readily parallelizable on GPUs, their impact on the overall wall clock time remains minimal.

Although mLSTM supports parallelization similar to approaches like FlashAttention~\citep{Dao:22,Dao:23} and GLA~\citep{Yang:23arxiv}, sLSTM cannot be parallelized efficiently because of its memory mixing arising from hidden-to-hidden interactions. Nonetheless, we have engineered an optimized CUDA version that incorporates GPU memory enhancements down to the register level, resulting in execution speeds typically within a factor of two compared to mLSTM.

\section{Related Work}
\label{sec:related}

\paragraph{Linear Attention.} A variety of techniques have been proposed to reduce the Transformer's quadratic complexity with respect to context length and render attention mechanisms linear in relation to sequence length. Synthesizer achieves this by generating synthetic attention weights that exclude direct interactions between tokens~\citep{Tay:20arxiv}. Linformer models self-attention using a low-rank decomposition and offers a linear approximation to the attention operation~\citep{Wang:20arxiv}. The Linear Transformer modifies the conventional attention calculation to yield a linear formulation \citep{Katharopoulos:20}. Performer approximates the attention softmax linearly through positive orthogonal random features~\citep{Choromanski:21}. Additionally, attention mechanisms have been replaced with rapid long-range convolutional operations, such as in Structured Global Convolution (SGConv)~\citep{Li:22} and the Hyena Hierarchy~\citep{Poli:23}.

\paragraph{State Space Models.} In recent years, State Space Models (SSMs) have gained significant traction due to their linear scaling with context length and their competitive results relative to Transformers. The introduction of the Structured State Space sequence model (S4)~\citep{Gu:21} marked an early milestone in this area, followed by a series of developments including the Diagonal State Space (DSS) model~\citep{Gupta:22}, Gated State Space (GSS) models~\citep{Mehta:22}, S5 model~\citep{Smith:22}, Bidirectional Gated SSM (BiGS)~\citep{Wang:22}, the H3 model~\citep{Fu:23}, and Mamba~\citep{Gu:24arxiv}.

\paragraph{Recurrent Neural Networks.} Recurrent Neural Networks (RNNs) have emerged as an alternative to Transformers and attention mechanisms, largely due to their ability to scale linearly with sequence length. Language modeling experiments utilizing Deep Linear Recurrent Units (LRUs) have yielded encouraging outcomes~\citep{Orvieto:23, De:24arxiv}, and similar gains have been observed with the Hierarchically Gated Linear RNN (HGRN)~\citep{Qin:23} and its successor HGRN2~\citep{Qin:24arxiv}. Among notable RNN-based architectures for large-scale language models, RWKV~\citep{Peng:23arxivshort,Peng:24arxivshort} has demonstrated strong results that are on par with those achieved by Transformers.

\paragraph{Gating.}
Gating is a central concept in LSTM architectures and has been independently developed and reinterpreted in numerous contemporary models. Variants employing gating mechanisms include HGRN~\citep{Qin:23}, HGRN2~\citep{Qin:24arxiv}, Gated Linear Attention (GLA)~\citep{Yang:23arxiv}, Gated State Space (GSS) models~\citep{Mehta:22}, Bidirectional Gated SSM (BiGS)~\citep{Wang:22}, Moving Average Equipped Gated Attention (MEGA)~\citep{Ma:22}, RWKV~\citep{Peng:23arxivshort}, as well as Mamba~\citep{Gu:24arxiv}.

\paragraph{Covariance Update Rule.} In order to improve storage capacity, we augmented the mLSTM cell with a matrix memory that employs a covariance update rule. This approach aligns with the update mechanisms utilized in methods such as Fast Weight Programmers \citep{Schmidhuber:92ncfastweights,Schlag:21}, RWKV-5 and RWKV-6~\citep{Peng:24arxivshort}, Retention~\citep{Sun:23arxiv}, Linear Transformer~\citep{Katharopoulos:20}, and HGRN2~\citep{Qin:24arxiv}.

\paragraph{Most Related.} The models most similar in principle to xLSTM are Retention~\citep{Sun:23arxiv}, RWKV~\citep{Peng:23arxivshort,Peng:24arxivshort}, and HGRN2~\citep{Qin:24arxiv}. These architectures all utilize the concepts of matrix memory and/or gating mechanisms. Nevertheless, unlike the recently introduced sLSTM, these models lack the capacity for memory mixing. The ability to mix memory is crucial for addressing state tracking challenges, making LSTMs inherently more expressive compared to State Space Models (SSMs) and Transformers~\citep{Merrill:24,Deletang:23}. State tracking plays an important role in tasks such as executing code or maintaining consistency of entities throughout lengthy texts.

\paragraph{Residually Stacking Architectures.} In alignment with the majority of modern large-scale deep learning frameworks, xLSTM models are designed by layering foundational modules via residual connections~\citep{Srivastava:15,He:16}.

Such an approach has facilitated the advancement of deep convolutional networks~\citep{He:16} as well as Transformers~\citep{Vaswani:17}. The adoption of Transformer architectures serves as the primary technology underpinning recent Large Language Models (LLMs) including GPT-3~\citep{Brown:20short}, ChatGPT~\citep{Schulman:22short}, GPT-4~\citep{Achiam:23gpt4short}, Megatron-LM~\citep{Shoeybi:19}, Gopher~\citep{Rae:21short}, ERNIE 3.0 Titan~\citep{Wang:21arxivshort}, GLaM~\citep{Du:21glamshort}, Chinese M6~\citep{Lin:21}, the multilingual AlexaTM 20B~\citep{Soltan:22}, OPT~\citep{Zhang:22opt}, Chinchilla~\citep{Hoffmann:22short}, BLOOM~\citep{Scao:22arxivshort}, GLM-130B~\citep{Zeng:22arxivshort}, LaMDA~\citep{Thoppilan:22short}, PaLM~\citep{Chowdhery:22short}, Llama~\citep{Touvron:23arxiv}, and Gemini~\citep{GeminiTeam:23arxiv,Reid:24arxiv}.

\section{Experiments}
\label{sec:Exp}

This section presents an empirical analysis of xLSTM, focusing on its performance in language modeling relative to other established techniques. Section~\ref{sec:ExpSyntethic} details xLSTM's effectiveness on a range of synthetic benchmarks. A comparison of the validation perplexities achieved by different language modeling approaches, all trained on 15B SlimPajama tokens~\citep{Soboleva:23}, is provided in Section~\ref{sec:ExpComparison}, alongside a series of ablation experiments on xLSTM. Furthermore, the scaling characteristics of each model are explored in the spirit of \citet{Kaplan:20} and \citet{Brown:20short}.

In Section~\ref{sec:ExpLanguage}, we present a more comprehensive language modeling assessment by training xLSTM and the top-performing models from Section~\ref{sec:ExpComparison} on 300B SlimPajama tokens~\citep{Soboleva:23}. Our analysis proceeds by evaluating: (1) each method’s ability to generalize to longer sequence lengths, (2) validation perplexity and performance on downstream tasks following~\citep{Sutawika:24}, (3) accuracy across the 571 text domains of the PALOMA benchmark~\citep{Magnusson:23arxivshort}, and (4) scaling trends in the presence of a twentyfold increase in training data.

Throughout all experiments, the format xLSTM[$a$:$b$] is adopted to indicate the proportion $a/b$ between mLSTM-based and sLSTM-based xLSTM blocks. For instance, xLSTM[7:1] specifies that among eight total blocks, seven utilize mLSTM and one utilizes sLSTM. When the aggregate number of blocks is 48, this notation corresponds to allocating 42 blocks to mLSTM and 6 blocks to sLSTM. Additionally, in every experiment, mLSTM blocks are preceded by up-projection modules, and sLSTM blocks are followed by post up-projection modules.

\subsection{Synthetic Tasks and Long Range Arena}
\label{sec:ExpSyntethic}
To begin, we examine how the novel exponential gating mechanism with memory mixing in xLSTM performs on formal language benchmarks~\citep{Deletang:23}. Next, we evaluate xLSTM's newly introduced matrix memory on the Multi-Query Associative Recall task~\citep{Arora:23arxiv}. Lastly, the capability of xLSTM to handle lengthy input sequences is measured using the Long Range Arena benchmark~\citep{Tay:21}.

\paragraph{Evaluation of xLSTM’s Exponential Gating Mechanism Incorporating Memory Mixing.} We evaluate the exponential gating with memory mixing feature introduced in xLSTM, which is hypothesized to facilitate state tracking problem resolution~\citep{Merrill:24, Merrill:23}. To this end, we adapt and broaden the suite of formal language benchmarks described by \citet{Deletang:23}, enhancing them to support multi-length training regimes aimed at probing extrapolation to longer sequences. Comprehensive descriptions of all tasks and a thorough presentation of empirical findings can be found in Appendix~\ref{sec:appExpSynthetic-formalLang}. 

We perform a systematic comparison between xLSTM and a range of alternative approaches including Transformer architectures, State Space Models, and standard Recurrent Neural Networks. Method performance is measured using token-level accuracy restricted to tokens pertinent to the respective task, normalized such that 0 corresponds to chance-level and 1 to perfect accuracy. Our analysis contrasts the capabilities of 2-block model architectures: xLSTM[0:1] (solely sLSTM), xLSTM[1:0] (solely mLSTM), xLSTM[1:1], alongside Llama, Mamba, RWKV, Retention, Hyena, LSTM, and LSTM employed within Transformer-style blocks (LSTM (Block)). Figure~\ref{fig:formal-main} visualizes the experimental outcomes.

Architectures lacking memory mixing—such as traditional Transformers and State Space Models—are incapable of addressing certain problems, including regular grammar recognition exemplified by the parity task. These empirical results corroborate prior work indicating the intrinsic limitations of Transformers and State Space Models relative to RNNs~\citep{Merrill:24, Merrill:23, Deletang:23}.

\paragraph{Test of xLSTM's Memory Capacities on Associative Recall Tasks.} This study investigates the effectiveness of the new matrix memory in xLSTM by evaluating its performance on the Multi-Query Associative Recall task~\citep{Arora:23arxiv}. In this task, each input sequence consists of key-value pairs, randomly sampled from an extensive vocabulary, that must be retained for future access. To increase the complexity beyond the original setup, we expand both the number of key-value pairs—up to 256—and the total context window—up to 2048 tokens. This adjustment provides a more comprehensive assessment of various models’ memory capacities. We benchmark 2-block model variants: Llama, Mamba, RWKV-5, RWKV-6, xLSTM[1:1], and xLSTM[1:0]. The criterion for evaluation is the models' accuracy in retrieving the correct key-value pairs. Transformers, like Llama, are known to possess memory capabilities scaling exponentially with coding dimension~\citep{Ramsauer:21}, making them a reference point for this benchmark. The comparative results are summarized in Figure~\ref{fig:mqar-main}.

xLSTM[1:1] demonstrates the strongest performance among all non-Transformer competitors, retaining this advantage even at smaller model sizes. Notably, the integration of the sLSTM block does not impair memory retention; rather, it appears to enhance the model's recall ability, particularly evident when tasked with 256 key-value pairs—the most challenging configuration. Further experimental findings are discussed in Appendix~\ref{sec:appExpSynthetic-mqar}, where analysis suggests that the improved memory capacity of xLSTM facilitates generalization to longer contexts than those encountered during training.

\paragraph{Test of xLSTM's Long Context Capabilities on Long Range Arena.} To evaluate how well xLSTM manages extended sequences and substantial contexts, we conduct a comparison among various approaches using the Long Range Arena~\citep{Tay:21}. The results show that xLSTM achieves robust and reliable outcomes across the full set of tasks, indicating that the architecture is highly adept at addressing multiple dimensions of long context challenges.

For more details, see Appendix~\ref{sec:appExpSynthetic-lra}.

\subsection{Method Comparison and Ablation Study}
\label{sec:ExpComparison}

In order to investigate the primary focus of this work—namely, evaluating the capabilities of our proposed LSTM variants under large-scale language modeling—we conduct experiments by training xLSTMs alongside Transformers, State Space Models, and additional baselines. Each model is trained on a dataset comprising 15B tokens from SlimPajama, employing a consistent auto-regressive framework. We evaluate these trained models on a shared validation set and conduct comprehensive ablation analyses to assess the behavior of the xLSTMs.

\paragraph{Benchmarking xLSTM Against Alternative Techniques.} To facilitate comparison, all models were trained on a 15B-token subset from the SlimPajama dataset~\citep{Soboleva:23} and evaluated by measuring perplexity on a held-out validation set. The range of methods under evaluation comprises: xLSTM (our proposed architecture), GPT-3 (Transformer)~\citep{Brown:20short}, Llama (Transformer)~\citep{Touvron:23arxiv}, H3 (State Space Model)~\citep{Fu:23}, Mamba (State Space Model)~\citep{Gu:24arxiv}, RWKV-4 (RNN)~\citep{Peng:23arxivshort}, RWKV-5 (RNN)~\citep{Peng:24arxivshort}, RWKV-6 (RNN)~\citep{Peng:24arxivshort}, GLA (linear Transformer)~\citep{Yang:23arxiv}, HGRN (RNN)~\citep{Qin:23}, HGRN2 (RNN)~\citep{Qin:24arxiv}, as well as RetNet (linear Transformer)~\citep{Sun:23arxiv}, Hyena (linear Transformer)~\citep{Poli:23}, xLSTM[1:0], and xLSTM[7:1] (refer to Section~\ref{sec:xlstm_blocks_notation} for clarification). 
All models were trained using mixed-precision, although RWKV-5, RWKV-6, GLA, and HGRN2 did not leverage PyTorch's automated mixed-precision training routines (see Appendix Section~\ref{sec:appGenTrainProc} for further discussion).
For the purpose of analysis, we group the architectures as follows: (a) Transformers, (b) State Space Models (SSMs), and (c) a unified category consisting of Recurrent Neural Networks (RNNs) and linear Transformers (the latter adopt a linear alternative to the attention computation in standard Transformers). Each implementation targets parity with GPT-3’s 350M parameter footprint, employing an embedding size of 1024 and utilizing 24 residual blocks. Note, GPT-3 employs shared parameters for token and output embeddings, thus has a marginally reduced overall parameter count. 
As shown in Table~\ref{tab:spaj15b_model_results}, xLSTM achieves superior validation perplexity relative to all other methods assessed. Additional experimental details are presented in Appendix~\ref{sec:appExpComparison}. Figure~\ref{fig:temp_sclaw15B} further illustrates the scaling properties uncovered in these experiments, implying continued robust performance of xLSTM with increasing model size.

\paragraph{Ablation Studies.}
Table~\ref{tab:spaj15b_model_results} along with Figure~\ref{fig:temp_sclaw15B} show that xLSTM performs strongly in language modeling tasks when trained on 15B tokens sourced from SlimPajama. To systematically isolate the impact of each architectural modification leading from LSTM to xLSTM, we incrementally transform a standard LSTM model into xLSTM. The process begins by placing LSTM layers within pre-LayerNorm residual backbone structures. The next modification involves expanding these layers into a post up-projection configuration. In the final stage, we introduce exponential gating alongside matrix memory mechanisms.

The results are shown in Table~\ref{tab:ablstudies} (top). 

Our ablation experiments indicate that the notable performance gains result from both the exponential gating mechanism and the matrix memory. Furthermore, recognizing the significant role gating plays in RNNs and State Space Models, we systematically explore various gating strategies. Table~\ref{tab:ablstudies} (bottom) reveals that configuring each gate to be trainable and responsive to input consistently enhances performance. Supplementary analyses of the distinct backbone modules can be found in Appendix~\ref{sec:appAblDetails}.

\subsection{xLSTM as Large Language Model}
\label{sec:ExpLanguage}

This investigation concludes with extensive language modeling experiments, evaluating the viability of xLSTM in the context of LLMs. To this end, we scale up the dataset and train using 300B tokens sourced from SlimPajama. Notably, this quantity of training data matches that employed by models such as Mamba~\citep{Gu:24arxiv} and Griffin~\citep{De:24arxiv}.

Our evaluation includes a comparison between xLSTM and representative models from the RWKV-4, Llama, and Mamba families—one from each methodological category discussed in Section~\ref{sec:ExpComparison}. RWKV-4 was chosen as the representative RNN model because suitable training precision configurations for RWKV-5, RWKV-6, and HGRN2 were identified only after the commencement of experiments involving 300B tokens (further details can be found in Appendix~\ref{sec:appGenTrainProc}). We train multiple model variants with parameter counts of 125M, 350M, 760M, and 1.3B, systematically probing their ability to extrapolate to longer sequence lengths, as well as evaluating validation performance. Additional analyses include assessments of downstream task proficiency, language modeling effectiveness on the 571 textual domains featured in the PALOMA benchmark, and an exploration of scaling laws exhibited by these models.

\paragraph{Sequence Length Extrapolation.}
\label{sec:spaj300B_ctx_extrapolation}
We begin by evaluating how well large-scale 1.3B parameter models—including xLSTM, RWKV-4, Llama, and Mamba—generalize to longer input sequences. While all models are trained with a maximum context length of 2048, their performances are assessed on inputs extending to 16384 tokens. The outcomes are illustrated in Figure~\ref{fig:spaj300B_seq_extrapolation}. Notably, among these approaches, xLSTM consistently achieves lower perplexity at increased context lengths.

\paragraph{Validation Perplexity and Downstream Tasks.}
\label{sec:spaj_downstream_eval}
Next, across every model scale, we assess xLSTM, RWKV-4, Llama, and Mamba by examining their validation perplexity on the SlimPajama dataset for next-token prediction, as well as their capabilities on downstream benchmarks focused on common sense reasoning.

In Table~\ref{tab:lmeval}, the third column presents the validation set perplexity results for various approaches. The xLSTM[1:0] and xLSTM[7:1] configurations consistently achieve the lowest perplexity scores across all model scales. The remaining columns in Table~\ref{tab:lmeval} summarize the outcomes on a range of downstream tasks. Across nearly all tasks and model sizes, xLSTM outperforms the alternatives, with the exception of certain scenarios in the ARC task, where Mamba occasionally leads. Further information can be found in Appendix~\ref{sec:appExpLanguage}.

\paragraph{Performance on PALOMA Language Tasks.} Thirdly, we assess the next-token prediction capabilities of xLSTM, RWKV-4, Llama, and Mamba across all model sizes using the PALOMA language tasks~\citep{Magnusson:23arxivshort}. Evaluation is based on perplexity scores for predicting the next token over 571 diverse text domains, encompassing sources from nytimes.com to r/depression on Reddit. Table~\ref{tab:ppleval} presents results, with perplexity metrics organized by general language modeling tasks (first seven columns) and more granular domain-specific benchmarks (last five columns). Notably, xLSTM[1:0] outperforms xLSTM[7:1] in these tasks.

xLSTM[1:0] achieves lower perplexity compared to Mamba on 568 of 571 text domains (99.5\%), outperforms Llama in 486 of 571 cases (85.1\%), and demonstrates lower perplexity than RWKV-4 on 570 of 571 domains (99.8\%). Additional information is available in Appendix~\ref{sec:appExpLanguage}.

\paragraph{Scaling Laws.} Fourth, we investigate the power-law scaling properties, which facilitate forecasting model performance at increased parameter scales~\citep{Kaplan:20,Brown:20short}. 
As depicted in Figure~\ref{fig:temp_sclaw300B}, all evaluated models exhibit comparable scaling trends, differing primarily in baseline performance (offsets). Among the models, RWKV-4 shows the lowest performance, with Llama and Mamba performing somewhat better. xLSTM outperforms Mamba, and the difference between xLSTM and Mamba mirrors the performance gap seen between Mamba and Llama. These scaling patterns suggest that as models become larger, xLSTM is likely to maintain or strengthen its advantage over both Transformer-based and State-Space model architectures.

\textbf{Generation Times and Maximal Throughput.} In our final set of evaluations, depicted in Figure~\ref{fig:inference_time_generation}, we assess the duration required for text generation, alongside peak throughput measurements presented in Figure~\ref{fig:inference_time_decoding_throughput} (left), for various xLSTM models at the 1.3B parameter level. For a comparative baseline, we use implementations of Mamba, Llama, and RWKV of comparable scale from HuggingFace, ensuring a static key-value cache is included for the Llama setup. At the time these tests were conducted, it was not possible to compile the Transformer Model with a full cache, nor compile the Mamba architecture using \texttt{torch.compile}. During generation tasks, all models employ a batch size of 1 and a pre-fill value of 16—a configuration that is particularly advantageous for Transformer-based architectures.

As depicted in Figure~\ref{fig:inference_time_generation}, both xLSTM and other recurrent models, specifically Mamba and RWKV-4, exhibit linear computational scaling, contrasting with the quadratic scaling behavior of Llama. For throughput assessment, we vary batch size and pre-fill specifically for the Llama model to gauge differences.

In Figure~\ref{fig:inference_time_decoding_throughput} (right), results indicate that, owing to its constant memory requirement, xLSTM supports substantially larger batch sizes than Llama, thereby achieving superior maximal throughput.

\section{Limitations}

(i) Unlike mLSTM, the sLSTM’s mechanism for mixing memory inherently prevents the use of operations that can be parallelized, making a rapid parallel execution infeasible. Despite this, we implemented a custom CUDA kernel for sLSTM, achieving performance that is currently less than double the runtime of our parallel mLSTM setup. (ii) The current mLSTM CUDA kernels have not been finely tuned, resulting in an implementation approximately four times slower than either FlashAttention or the scan procedure utilized in Mamba. Significant speed-ups might be achievable by applying kernel optimizations similar to those in FlashAttention. (iii) Managing the matrix memory in mLSTM entails substantial computational cost, as it requires the manipulation of $d \times d$ matrices. However, because both memory updates and retrievals are parameter-free and lend themselves to standard matrix-parallel operations, the added computation from the matrix structure imposes minimal runtime overhead. (iv) Careful selection of the initial values for forget gates is important. (v) While the size of the matrix memory in mLSTM does not scale with sequence length, lengthening sequences can saturate the memory’s capacity, especially when dealing with very long contexts. Despite this, for contexts as long as 16k, this has not proven restrictive (see Section~\ref{sec:spaj300B_ctx_extrapolation}). (vi) Owing to the intensive computational requirements of large-scale language model experiments, we did not thoroughly fine-tune either the architecture or the hyperparameters, particularly in the larger xLSTM variants. Unlocking the full capability of xLSTM will likely require an exhaustive process of architectural and hyperparameter optimization.

\section{Conclusion}
We have provided a partial answer to our initial query: What is achievable in language modeling when scaling LSTM architectures to the billion-parameter regime? Our findings suggest that these models can perform on par with existing technologies, such as Transformers and State Space Models. By introducing the xLSTM variant—which incorporates exponential gating mechanisms featuring memory mixing and a revised memory architecture—we observe that xLSTM achieves competitive results in language modeling benchmarks against state-of-the-art models, including Transformers and State Space Models. The observed scaling behavior implies that larger xLSTM models have the potential to rival present-day Large Language Models predominantly based on the Transformer framework. Furthermore, xLSTM may also substantially benefit applications in other domains, such as Reinforcement Learning, Time Series Forecasting, and the simulation of physical systems.

\end{document}